\section{The Decision Procedures}
\label{sec:decision}

We now assemble the decision procedures for disjunctive entailment: over
arbitrary trees, where the previous sections apply directly; over regular
trees, where the answers are the same; over finite trees, where they differ
only through satisfiability; and for the signature $\BF$, which reduces to
$\BFT$. \Cref{thm:main} then follows with \cref{prop:sentences}. Proofs
omitted here are in \cref{app:classes}.

\subsection{Arbitrary Trees}\label{sec:decision-arbitrary}

\begin{theorem}[Arbitrary trees]\label{thm:decide-arbitrary}
Disjunctive entailment between flat constraints over $\BFT$ is decidable
over $\Trees$, uniformly in the number of queries, the arity and the
variance.
\end{theorem}

\begin{proof}
Let $\Gamma$ be a constraint and $(x_1,y_1),\dots,(x_k,y_k)$ queries.
\begin{enumerate}
\item Decide whether $\Gamma$ is satisfiable over $\Trees$
  (\cref{lem:sat-decidable}). If it is not, the entailment holds.
\item Otherwise, for each of the $2^k$ mode vectors $\bar\mu$, compute the
  reader $M_{\bar\mu}$ of \cref{def:compile} and test all tuples of paths of
  total length at most $\Hole(k,N)$, where $N$ is the number of its states,
  for being holes. The entailment fails if and only if a hole is found.
  For $k=0$ there is one mode vector, the empty one, and $\Hole(0,N)=0$; the
  empty tuple is a hole, and the entailment fails, as it must for a
  satisfiable constraint.
\end{enumerate}
The second step is correct by \cref{thm:reduction,thm:small-hole}. Whether a tuple is a hole is checked directly from
\cref{def:hole}.
\end{proof}

\begin{lemma}[Satisfiability]\label{lem:sat-decidable}
It is decidable whether a flat constraint over $\BFT$ is satisfiable over
$\Trees$.
\end{lemma}

Satisfiability is in fact decidable in polynomial time
\cite{KozenPalsbergSchwartzbach94,PalsbergWandOKeefe97}. In our setting it
is the question whether the least closure is consistent and contains no
forbidden atom (\cref{sec:readers-atoms}), and for each signed variable the
paths of its atoms in the least closure form a regular language accepted by
the automaton of lower runs (\cref{app:reduction-lc}).

\subsection{Regular Trees}\label{sec:decision-regular}

\begin{theorem}[Regular trees]\label{thm:decide-regular}
For every flat constraint $\Gamma$ over $\BFT$ and all queries,
\[
  \Gamma\entails_{\TreesReg}x_1\le y_1\vee\dots\vee x_k\le y_k
  \iff
  \Gamma\entails_{\Trees}x_1\le y_1\vee\dots\vee x_k\le y_k .
\]
\end{theorem}

\begin{proof}[Proof sketch]
A regular countermodel is a countermodel, which gives the implication from
right to left. Conversely, let $\rho$ be a countermodel over $\Trees$, with
seeds $A_i$ and exclusions $B_i$ as in \cref{sec:readers-modes}. By
\cref{thm:realization} the hull of the seeds is a solution set that
contains the seeds and none of the exclusions, so the corresponding
assignment is a countermodel. For each signed variable $x^s$, the paths $w$
with $(x^s,w)$ in the hull form a regular language: those from the least
closure do, and by \cref{lem:cuts} those derived from the seed $(p_i,w_i)$
are the words $v\,\wsuf{w_i}a\,\sigma$, for a cut $a$ of $w_i$, a word $v$
with $\Dn(x^s,v)\cap\Up(p_i,\wpre{w_i}a)\ne\varnothing$ and an arbitrary
word $\sigma$. A tree whose two cones are regular languages has finitely
many subtrees, because the subtree at a position is determined by the
residuals of the two cones and the polarity of the position
(\cref{lem:regular-cones}).
\end{proof}

So the hull of the seeds, the least countermodel with the given witnesses,
is a regular tree in every component, and it is computable from a hole.

\subsection{Finite Trees}\label{sec:decision-finite}

Over finite trees a constraint can be unsatisfiable although it is
satisfiable over $\Trees$; for a unary constructor, $x\le f(x)$, $f(x)\le x$
is an example. This is
the only difference.

\begin{theorem}[Finite trees]\label{thm:decide-finite}
For every flat constraint $\Gamma$ over $\BFT$ and all queries,
$\Gamma\entails_{\TreesFin}x_1\le y_1\vee\dots\vee x_k\le y_k$ holds if and
only if $\Gamma$ is unsatisfiable over $\TreesFin$ or
$\Gamma\entails_{\Trees}x_1\le y_1\vee\dots\vee x_k\le y_k$. Satisfiability
over $\TreesFin$ is decidable, and hence so is disjunctive entailment over
$\TreesFin$.
\end{theorem}

The nontrivial direction is to turn a countermodel over $\Trees$ into a
finite one when $\Gamma$ has some finite solution. A countermodel is
described by finitely many \emph{observations}: the $k$ seeds belong to its
set of atoms and the $k$ exclusions do not. Each observation is equivalent
to a bound of a variable by a finite tree. For instance, for a covariant
constructor, $(x^{\sgp},w)\in S_\rho$ says that $\rho(x)$ lies above the
finite tree that has $f$ along the path $w$, $\Top$ at its end and $\Bot$
beside it, and $(x^{\sgp},w)\notin S_\rho$ says that $\rho(x)$ lies below the
finite tree that has $f$ along $w$ and at its end, and $\Top$ everywhere
else (\cref{lem:observations}). The theorem then follows from a pinning
property of finite solutions:

\begin{lemma}[Pinning]\label{lem:pinning}
Suppose that $\Gamma$ has a solution over $\TreesFin$ and that some solution
of $\Gamma$ over $\Trees$ satisfies finitely many bounds $g\le x$ and $x\le
g'$ of variables by finite trees. Then some solution of $\Gamma$ over
$\TreesFin$ satisfies these bounds.
\end{lemma}

Its proof (\cref{app:classes-finite}) first treats a covariant
constructor. It compiles the bounds into flat literals on fresh variables,
which only hang finite trees below or above existing variables, and shows
that they create no \emph{cycle clash}, the obstruction to finite
satisfiability of \cite{Zhao26}. With variance, it translates the
constraint into a covariant constraint on the signed variables and makes the
finite solution obtained there symmetric under the exchange of the two signs,
by a median of trees.

\subsection{The Signature \texorpdfstring{$\BF$}{\{bot, f\}}}\label{sec:decision-bf}

Let $f$ be covariant. The trees over $\BF$ are the trees over $\BFT$ without
$\Top$, with the same order, so a constraint over $\BF$ can be read over
$\BFT$; but it may then have additional solutions, in which $\Top$ occurs. A
single \emph{guard} variable excludes them.

\begin{lemma}[Guard]\label{lem:guard}
Let $f$ be covariant. If $u\le f(u,\dots,u)$, then $\Top$ occurs neither in
$u$ nor in any tree $t\le u$.
\end{lemma}

\begin{proof}
Suppose $\Top$ occurs in $u$, and let $\pi$ be a position of minimal length
with $u(\pi)=\Top$. Since $u\le f(u,\dots,u)$ and $u\ne\Top$, either
$u=\Bot$, which is impossible, or $u=f(u_1,\dots,u_n)$ with $u_i\le u$ for
every $i$. Then $\pi=i\pi'$ and $u_i(\pi')=\Top$. Let $v$ be the longest
prefix of $\pi'$ that is a position of $u$. If $v=\pi'$, then $u_i\le u$
gives $u(\pi')=\Top$; otherwise $u(v)$ is a leaf above the label $f$ of
$u_i(v)$, so $u(v)=\Top$. In both cases $\Top$ occurs in $u$ at a position
shorter than $\pi$. The same argument, applied to a position of $\Top$ in a
tree $t\le u$, shows that $\Top$ occurs in $u$.
\end{proof}

For a constraint $\Gamma$ over $\BF$ with variables $V$, let $\Gamma^{g}$ be
$\Gamma$ together with the literals $u\le f(u,\dots,u)$ and $x\le u$ for
every $x\in V$, where $u$ is a fresh variable.

\begin{theorem}[Reduction of $\BF$]\label{thm:decide-bf}
Let $f$ be covariant and $\Cls$ one of $\Trees$, $\TreesReg$, $\TreesFin$.
A constraint $\Gamma$ over $\BF$ entails $x_1\le y_1\vee\dots\vee x_k\le
y_k$ over the trees of $\Cls$ without $\Top$ if and only if $\Gamma^{g}$
entails it over $\Cls$ as a constraint over $\BFT$.
\end{theorem}

\begin{proof}
By \cref{lem:guard}, the restriction of a solution of $\Gamma^{g}$ to $V$ is
a solution of $\Gamma$ without $\Top$. Conversely, a solution of $\Gamma$ in
trees of $\Cls$ without $\Top$ extends to a solution of $\Gamma^{g}$ in
$\Cls$: over $\Trees$ and $\TreesReg$ let $u$ be the tree $\fall$ with $f$
at every path (for $n=0$, the tree $f$), which lies above every tree without
$\Top$ and satisfies $u=f(u,\dots,u)$; over $\TreesFin$ let $u$ be the
complete tree of a depth $d$ that exceeds the lengths of all positions of
the finitely many values, with $f$ at all paths shorter than $d$ and $\Bot$
at the paths of length $d$. The queries do not involve $u$, so the two
entailments have the same countermodels up to the value of $u$.
\end{proof}

The guard is needed: at arity $0$, the inequality $x\le f$ holds for both
trees $\Bot$ and $f$ over $\BF$, but it fails for $x=\Top$.

\subsection{Proof of the Main Theorem}\label{sec:decision-main}

\begin{proof}[Proof of \cref{thm:main}]
Disjunctive entailment between flat constraints over $\BFT$ is decidable
over $\Trees$ by \cref{thm:decide-arbitrary}, over $\TreesReg$ by
\cref{thm:decide-regular}, and over $\TreesFin$ by
\cref{thm:decide-finite}. Each of the three classes contains $\Bot$ and
$\Top$ and is closed under the constructor, so \cref{prop:sentences} decides
the Boolean combinations of $\Sigma_1$ and $\Pi_1$ sentences over each
class, and its second statement together with \cref{thm:decide-regular}
shows that truth over $\TreesReg$ and over $\Trees$ agree. For $\BF$ with a
covariant constructor, the trees without $\Top$ of each class contain $\Bot$
and are closed under the constructor, \cref{lem:clause} holds for terms
without $\Top$, and \cref{thm:decide-bf} decides the resulting disjunctive
entailments; regular and arbitrary trees agree again, because they agree for
the guarded constraints. All steps are uniform in the arity and the
variance.
\end{proof}
